(1) 回路满足 $V_0=RI+q/C$、$I=\dot q$。设时间常数 $\tau=RC$，线性一阶方程解为 $q(t)=CV_0(1-e^{-t/\tau})$、$I(t)=(V_0/R)e^{-t/\tau}$，其中 $q(0)=0$，$q(\infty)=CV_0$，符合充电终态电流为零。
(2) 电源提供总功 $W_s=\int_0^\infty V_0I(t)\,dt=V_0[q(\infty)-q(0)]=CV_0^2$。电容储能为 $U_C=q(\infty)^2/(2C)=CV_0^2/2$。因此电容只储存电源所供能量的一半，不能忽略串联电阻消耗。
(3) 焦耳热率 $P_R=I^2R=(V_0^2/R)e^{-2t/(RC)}$，积分为 $Q_R=\int_0^\infty P_Rdt=CV_0^2/2$。它与电容储能之和是 $CV_0^2$，恰好等于电源总功，独立核对了全过程的热力学第一定律。
(4) 放电阶段 $q(t)=CV_0e^{-t/(RC)}$，按充电正向定义电流为 $I(t)=-V_0e^{-t/(RC)}/R$。总电阻热量仍为 $\int_0^\infty I^2Rdt=CV_0^2/2$。改变 $R$ 只改变时间尺度 $RC$ 与瞬时电流，不改变整个过程的耗散总能量。
